\clearpage
\pdfbookmark[1]{Kugelfunktionen}{spherical-harmonics}
\begin{center}\textbf{Kugelfunktionen}\end{center}
\noindent Diese mit KI erstellte Fassung in modernem Deutsch beruht auf dem deutschen Original \textit{Kugelfunctionen} in Carl Friedrich Gauss, \textit{Werke}, Band~V, S.~630--631, herausgegeben von der Königlichen Gesellschaft der Wissenschaften zu Göttingen. Siehe das \href{https://gdz.sub.uni-goettingen.de/id/PPN236006339}{Göttinger Digitalisat} und den zweiten Abdruck (1877), \href{https://archive.org/download/Werkecarlf05gausrich/page/n641}{S.~630} und \href{https://archive.org/download/Werkecarlf05gausrich/page/n642}{S.~631}. Die Abbildung ist in bearbeitbarem TikZ gezeichnet.
\par\smallskip\noindent Zur Terminologie siehe T. M. Dunster, \href{https://dlmf.nist.gov/14.30.iv}{NIST Digital Library of Mathematical Functions, §14.30(iv)}, Version 1.2.8 vom 15. September 2026, und Jens Wirth, \href{https://pnp.mathematik.uni-stuttgart.de/iadm/Wirth/Skripte/HarmonischeAnalysis/}{\textit{Harmonische Analysis}}, §4.4, Beispiele 4.27--4.28, Universität Stuttgart.
\clearpage
